% Generated draft: requires mathematical and source review.
\documentclass[11pt]{article}
\usepackage[margin=25mm]{geometry}
\usepackage{fontspec}
\setmainfont{Latin Modern Roman}
\usepackage{amsmath,amssymb,mathtools,unicode-math}
\setmathfont{Latin Modern Math}
\usepackage{enumitem,needspace,xurl,hyperref}
\setlength{\parindent}{0pt}
\setlength{\parskip}{5pt}
\setlength{\emergencystretch}{3em}
\newcommand{\N}{\mathbb{N}}
\newcommand{\Z}{\mathbb{Z}}
\newcommand{\Q}{\mathbb{Q}}
\newcommand{\R}{\mathbb{R}}
\newcommand{\C}{\mathbb{C}}
\newcommand{\F}{\mathbb{F}}
\newcommand{\A}{\mathbb{A}}
\newcommand{\PP}{\mathbb{P}}
\newcommand{\E}{\mathbb{E}}
\newcommand{\eps}{\varepsilon}
\newcommand{\dd}{\,\mathrm{d}}
\begin{document}
\section*{UnsolvedMath ID: 30002121}
\textbf{Source problem code:} catalogue entry
\subsection{Definitions and mathematical statement}
A convex \(d\)-polytope is a bounded convex hull of finitely many points in \(\R^d\) whose affine hull is \(\R^d\). Two such polytopes have the same combinatorial type if their face lattices, ordered by inclusion, are isomorphic. They are projectively equivalent if a projective transformation of \(\mathbb{RP}^d\) maps one onto the other, with both polytopes regarded as subsets of the standard affine chart. A polytope \(P\) is projectively unique if every convex \(d\)-polytope having the same combinatorial type as \(P\) is projectively equivalent to \(P\).

For each integer \(d\), let \(\mathcal U_d\) denote the set of combinatorial types of projectively unique convex \(d\)-polytopes. The supplied clean statement asks to determine the set
\[
D_{\infty}=\bigl\{d\in\Z:4\le d\le80\ \text{and}\ \mathcal U_d\ \text{is infinite}\bigr\}.
\]
Equivalently, determine exactly which integers \(d\) with \(4\le d\le80\) satisfy the following condition: for every positive integer \(N\), there exist projectively unique convex \(d\)-polytopes \(P_1,\ldots,P_N\) whose face lattices are pairwise nonisomorphic. The vertex counts are unrestricted.

The original statement additionally asks, for fixed dimension \(d\), whether \(\mathcal U_d\) is finite and how the dimension of the realization space of a \(d\)-polytope grows with its size. That additional growth question is not specified precisely enough to formalize uniquely: the record does not define size, the realization-space convention, or the intended growth quantity and quantifiers. In particular, it does not specify whether realizations are taken modulo affine or projective transformations, or whether the question concerns bounds or extrema over polytopes of a given size.

The supplied qualifications require separate review. The statement-verification field says that the realization-space question is answered in the cited article and that finiteness fails in every dimension at least \(81\). The supplied dated literature assessment instead says that the full growth question remains unresolved and cites a construction in dimension \(69\), which lies inside the clean statement's range. Thus the interval \(4\le d\le80\) is retained as the requested classification range, without asserting that every dimension in it is unresolved. These source claims are recorded as supplied qualifications, not as independently verified mathematical conclusions.
\subsection{Short English statement}
Determine which integer dimensions \(4\le d\le80\) admit infinitely many combinatorial types of convex polytopes whose realizations are unique up to projective transformation. The original record also asks an insufficiently specified realization-space growth question. Its supplied literature and verification notes disagree about the remaining scope, so the claimed unresolved range and status require human review.
\subsection{Sources}
\begin{itemize}
\item Supplied local UnsolvedMath registry record, permanent ID \(30002121\), exact source problem code \texttt{catalogue entry}, titled ``Projectively Unique Polytopes in Intermediate Dimensions,'' in the Oberwolfach Reports Open Problems collection. The record gives proposed year \(2012\) and names no proposer.
\item Supplied report citation: \emph{Discrete Differential Geometry (2013)}, Oberwolfach Reports, DOI \texttt{10.4171/owr/2012/34}, pp.\ 2077--2137. Supplied source link: \url{https://doi.org/10.4171/owr/2012/34}. The background describes Oberwolfach Reports as documenting mathematical workshops and their open-problem discussions. The citation's year \(2013\), the DOI's \(2012\) component, and the proposed year are retained without reconciliation.
\item Supplied literature source labeled ``Adiprasito Ziegler'': \url{https://arxiv.org/abs/1212.5812}. No further bibliographic identification is supplied in the record.
\item Supplied literature source labeled ``OWR report'': \url{https://ems.press/content/serial-article-files/46402}.
\item Literature assessment dated exactly \texttt{2026-08-21}: supplied dated curation, not independently verified. Its status is ``partially solved.'' It says that the finiteness question is positive in dimensions two and three and negative in sufficiently high fixed dimensions, explicitly dimension sixty nine in the primary construction; intermediate dimensions and the full growth question remain unresolved. The background describes this as literature triage derived from supplied OWR curation and requests independent verification of time-sensitive status and cited claims.
\item The record separately labels the clean statement \texttt{corrected\_verified}. Its supplied verification note says that the report's second realization-space problem is answered in the same article, that its first problem is answered negatively for every dimension at least \(81\), and that the clean statement isolates the unresolved intermediate-dimensional range. This label and explanation are source assertions, not independent verification by this draft. The registry record was updated on \texttt{2026-09-16T12:06:53Z}.
\end{itemize}
\end{document}
